\documentclass[11pt,a4paper]{article}
\input{glyphtounicode}
\pdfgentounicode=1
\usepackage{cmap}
\usepackage[margin=0.72in]{geometry}
\usepackage[T1]{fontenc}
\usepackage[utf8]{inputenc}
\usepackage[hidelinks]{hyperref}
\usepackage{array}
\usepackage{xcolor}

\definecolor{accent}{HTML}{B45309}
\definecolor{textgray}{HTML}{334155}

\setlength{\parindent}{0pt}
\setlength{\parskip}{5pt}
\renewcommand{\familydefault}{\sfdefault}

\newcommand{\cvsection}[1]{%
  \vspace{8pt}
  {\large\bfseries\color{accent} #1}\\[-2pt]
  \rule{\linewidth}{0.4pt}
  \vspace{2pt}
}

\newcommand{\entry}[4]{%
  \textbf{#1}\hfill{\small #2}\\
  \textit{#3}\\
  #4
  \vspace{4pt}
}

\begin{document}

{\LARGE\bfseries Syed Tahir Ali Shah}\\
{\large Ph.D. Researcher in Mechanical Engineering, Embedded IoT, Affective Computing, and Smart Mobility}\\[4pt]
\href{mailto:syedtahiralishah@ua.pt}{syedtahiralishah@ua.pt} \quad
\href{mailto:syedtahirkzm1@gmail.com}{syedtahirkzm1@gmail.com}\\
Centre for Mechanical Technology and Automation (TEMA), Department of Mechanical Engineering, University of Aveiro, Portugal\\
\href{https://orcid.org/0000-0003-0001-1810}{ORCID} \quad
\href{https://scholar.google.com/citations?user=VED2-h4AAAAJ}{Google Scholar} \quad
\href{https://www.cienciavitae.pt//F712-C83E-0EFB}{Ciencia Vitae} \quad
\href{https://pt.linkedin.com/in/syed-tahir-ali-shah-a1582897}{LinkedIn}

\cvsection{Profile}
Multidisciplinary Ph.D. researcher working across Mechanical Engineering, Computer Science, embedded Internet-of-Things systems, wearable sensing, edge intelligence, and human-centred smart mobility. Current research focuses on emotion-aware and physiology-aware mobility systems that use wearable signals, embedded machine learning, and IoT integration to support safer and more adaptive mobility experiences.

\cvsection{Education and Research Appointments}
\entry{Ph.D. in Mechanical Engineering}{2023/11/01--2027/06/30 (ongoing)}{University of Aveiro, Portugal}{Focus: emotion-aware IoT-integrated mobility systems.}
\entry{Research Trainee}{From 2026/01/01}{Fundacao para a Ciencia e a Tecnologia, Portugal}{Host: University of Aveiro, Department of Mechanical Engineering.}
\entry{Scholarship Holder}{2023/11/01--2025/12/31}{University of Aveiro, Portugal}{Ph.D. research work in the AM2R and MariaBike research context.}
\entry{M.Sc. Computer Science}{2016/01/10--2018/08/31}{Iqra University - Islamabad Campus, Pakistan}{Classification: 14/20.}
\entry{B.Sc. Computer Science}{2011/06/01--2016/01/01}{Karakoram International University, Pakistan}{Major: software engineering, database systems, data structures, and algorithm design.}

\cvsection{Research Interests}
\begin{itemize}
  \item Affective computing and emotion-aware IoT systems.
  \item Embedded machine learning, TinyML, mobile runtimes, and edge intelligence.
  \item Wearable-driven physiological monitoring for mobility safety.
  \item Human-centred smart mobility, e-bikes, and intelligent transport systems.
\end{itemize}

\cvsection{Selected Publications}
\begin{itemize}
  \item Syed Tahir Ali Shah, J.P. Santos, Gabriel Constantinescu, Jose M. Fernandes, A.B. Pereira. ``Physio-Digital Twin for Human-Centered IoT Mobility: A Proof-of-Concept Implementation on the MariaBike Platform.'' ICICT 2026. Best Paper Award. DOI: \href{https://doi.org/10.1145/3803291.3803295}{10.1145/3803291.3803295}.
  \item Syed Tahir Ali Shah, J. M. Fernandes, J. P. Santos, G. Constantinescu, Antonio B. Pereira. ``The Micro-Mobility Sensing Gap: A Systematic Review of Physiological Safety Monitoring from Cycling to E-Scooters.'' Sensors, 2026. DOI: \href{https://doi.org/10.3390/s26041110}{10.3390/s26041110}.
  \item Gabriel Constantinescu, Syed Tahir Ali Shah, Jose Paulo Oliveira Santos, Joao Manuel Cardoso, Mario Jorge de Sousa Henriques, Antonio Manuel de Bastos Pereira. ``Tensile and Flexural Behavior of Biaxial Non-Crimp-Fabric Composites for Two-Wheeled Electric-Vehicle Chassis.'' Fibers, 2026. DOI: \href{https://doi.org/10.3390/fib14050061}{10.3390/fib14050061}.
  \item Syed Tahir Ali Shah. ``Extending IoT Networks with the Human Factor: An Example for an E-bike Scenario.'' 7th International Conference on Technologies for the Wellbeing and Sustainable Manufacturing Solutions, 2024. DOI: \href{https://doi.org/10.48528/mehv-z227}{10.48528/mehv-z227}.
\end{itemize}

\cvsection{Projects and Software}
\entry{MariaBike Dashboard}{2025}{Featured software}{Real-time e-bike dashboard with BLE and MQTT integration for physiological monitoring in mobility.}
\entry{FCT BDANA PhD Work Programme}{2026--2030}{Ph.D. studentship}{Embedded intelligent systems for personalized and sustainable mobility.}
\entry{AM2R Mobilizing Agenda}{2023--2025}{Ph.D. student fellow}{Business innovation in the two-wheel sector at the University of Aveiro.}

\cvsection{Service and Mentorship}
\entry{Scientific Mentorship (CHAI I)}{2024/10/01--2025/06/30}{University of Aveiro}{Mentoring student Joao Bispo in Scientific Mentorship I.}
\entry{Thesis Supervisor}{2022/06/01--2023/06/30}{University of Baltistan, Pakistan}{Topic: IoT-Based Smart Farming System with Real-Time Soil Moisture and Weather Monitoring.}
\entry{Thesis Co-supervisor}{2020/01/01--2021/01/31}{University of Baltistan, Pakistan}{Topic: Real-Time Emotion Recognition from Facial Expressions Using Deep Learning.}

\cvsection{Languages}
\begin{tabular}{>{\bfseries}p{0.24\linewidth}p{0.68\linewidth}}
Urdu & Native / mother tongue\\
English & Upper intermediate to proficient across communication skills\\
Portuguese & Beginner to intermediate working proficiency\\
\end{tabular}

\vfill
{\small This CV is generated for the personal academic website. For the complete official record, visit \href{https://www.cienciavitae.pt//F712-C83E-0EFB}{Ciencia Vitae}.}

\end{document}
